%% kapa: exam=1 round=30 year=2019 period=2 form=A subject=accounting subject_ko=회계학 seq=4 number=44
%% source: Q-Net 감정평가사 기출문제 (cst00309) artlSeq=5205205 file=extracted/1008/5205205/2019년 제30회 감정평가사 1차시험 문제지/2019년 제30회 감정평가사 1차시험 문제지/2019년도 제30회 감정평가사 1차 2교시 A형.hwp
%% conversion: kordoc 4.24.51 -> tools/gen1.py (auto); review=proofread-sonnet; flags=none
\begin{question}{44}
(주)감평은 20x1년 초 기계장치(취득원가 ￦6,000,000, 내용연수 5년, 잔존가치 ￦0)를 취득하여 정액법으로 감가상각하고 있다. 20x1년 말 이 기계장치에 손상징후가 존재하여 회수가능액을 추정한 결과 회수가능액이 ￦2,232,000으로 추정되었다. (주)감평은 동 금액과 장부금액 간의 차이가 중요한 것으로 판단하여 손상차손을 인식하였다. 한편, 20x2년 말 기계장치의 회수가능액이 ￦4,000,000으로 회복된 것으로 추정될 경우, (주)감평이 20x2년 말 인식할 손상차손환입액은? (단, 기계장치에 대하여 원가모형을 적용한다.)
\begin{choices}
  \item ￦1,574,000
  \item ￦1,926,000
  \item ￦2,138,000
  \item ￦2,326,000
  \item ￦2,568,000
\end{choices}
\end{question}
\answer{2}
